\section{System overview}\label{sec:overview}

\Cref{fig:architecture} shows the complete system. Stage~1
(\cref{sec:stage1}) produces the frozen preprocessing outputs on the left;
perception (\cref{sec:stage3}) turns them into frozen observations; the
memory organisations, retrieval and answering (\cref{sec:stage4}) and the
evaluation (\cref{sec:stage5}) follow.

\begin{figure}[ht]
\centering
\resizebox{\linewidth}{!}{%
\begin{tikzpicture}[
  font=\small,
  box/.style={draw, rounded corners=2pt, align=center, minimum height=0.9cm,
              inner sep=4pt, fill=#1},
  box/.default=white,
  arr/.style={-{Stealth[length=2mm]}, thick},
  node distance=0.55cm and 0.6cm]
  \node[box=gray!12] (video) {video\\\texttt{data/raw}};
  \node[box=blue!8, right=of video, yshift=0.7cm] (audio) {audio\\16 kHz mono};
  \node[box=blue!8, right=of audio] (asr) {ASR\\segments};
  \node[box=blue!8, right=of video, yshift=-0.7cm] (frames) {frames\\$N$ per chunk};
  \node[box=blue!8, right=of asr, yshift=-0.7cm] (chunks) {chunks\\$[u_i,v_i)$};
  \node[draw, dashed, rounded corners, inner sep=6pt,
        fit=(audio)(asr)(frames)(chunks),
        label={[font=\footnotesize]above:Stage 1 (frozen, hashed)}] (s1) {};
  \node[box=orange!12, right=of chunks] (perc) {perception $\phi$\\(VLM)};
  \node[box=orange!12, right=of perc] (obs) {$o_{1:n}$\\frozen};
  \node[box=green!10, below=1.2cm of obs, xshift=-2.6cm] (c1) {C1\\summary};
  \node[box=green!10, right=0.25cm of c1] (c2) {C2\\FIFO log};
  \node[box=green!10, right=0.25cm of c2] (c3) {C3\\structured};
  \node[box=green!10, right=0.25cm of c3] (c4) {C4\\unbounded};
  \node[box=purple!8, below=1.0cm of c2, xshift=0.8cm] (ans)
       {retrieval $R$ + answerer $A$\\$(\hat a, C)$, cited records};
  \node[box=gray!12, below=0.7cm of ans] (eval)
       {evaluation: accuracy, survival, grounding};
  \draw[arr] (video) |- (audio);
  \draw[arr] (video) |- (frames);
  \draw[arr] (audio) -- (asr);
  \draw[arr] (asr) -| (chunks);
  \draw[arr] (frames) -| (chunks);
  \draw[arr] (chunks) -- (perc);
  \draw[arr] (perc) -- (obs);
  \foreach \c in {c1,c2,c3,c4} \draw[arr] (obs.south) -- (\c.north);
  \foreach \c in {c1,c2,c3,c4} \draw[arr] (\c.south) -- (ans.north);
  \draw[arr] (ans) -- (eval);
\end{tikzpicture}}
\caption{Architecture. Every memory organisation C1--C4 reads the same frozen
observations and is queried by the same retrieval allowance and answerer.}
\label{fig:architecture}
\end{figure}

\subsection{Design invariants}

\paragraph{Memory-agnostic perception.}
The perception model describes each chunk without seeing the memory. Linking a
``red mug'' in one chunk to an entity observed twenty minutes earlier is the
job of the memory, not of perception. Conditioning perception on memory would
be a natural engineering choice, since it gives the model context, but it
would make the observations depend on the organisation under test, and the
comparison of \cref{tab:conditions} would no longer be controlled.

\paragraph{Event-sourced memory.}
Every operation (insert, merge, update, compress, delete) is appended to an
operation log, and the memory after chunk $i$ is obtained by replaying the log
up to $i$. This yields three things at no extra cost: the memory-over-time view
required by the demo, an audit trail for the failure attribution of
\cref{sec:attribution}, and a correctness check (replaying the full log
reproduces the live memory exactly; this was verified for every memory
organisation during development).

\paragraph{Provenance preservation.}
The invariant of \cref{sec:formulation}, that merging and compression take the
union of provenance sets, is enforced in code and was checked after every
operation type during development. Without it, grounding (RQ4) and fact survival (RQ2) could not
be measured.

\paragraph{Single budget unit.}
The budget is measured in tokens of the answerer's tokenizer applied to the
serialised memory, for every organisation, so ``equal budget'' has one
meaning.

\subsection{Implications for Stage 1}
These invariants put four requirements on preprocessing: (i)~chunk boundaries
must be fixed once, because provenance is expressed in chunk indices;
(ii)~timestamps must be accurate and verified, because grounding compares them
with ground truth; (iii)~the inputs to perception must be byte-identical
across all later runs; (iv)~any accidental modification must be detectable.
\Cref{sec:stage1} describes how each is met.
